% --- Capítulo 5: Entorno de desarrollo y herramientas tecnológicas ---
\chapter{Entorno de desarrollo y herramientas tecnológicas}

\section{Lenguaje de programación}
% Breve justificación del uso de Python 3.x como estándar en la industria 
% del Machine Learning y su rico ecosistema de librerías científicas.

\section{Framework principal de Deep Learning: PyTorch}
% Sustituimos TensorFlow por PyTorch, que es lo que realmente usas.
\subsection{Justificación de la elección de PyTorch}
% Mencionar su flexibilidad, el uso del grafo computacional dinámico 
% y su facilidad para depurar modelos personalizados (como tu HybridViTCNN).
\subsection{Ecosistema Torchvision}
% Explicar el uso de torchvision para cargar los pesos preentrenados 
% (ResNet18, ViT) y aplicar transformaciones de imágenes (Normalización, Resize).

\section{Herramientas de procesamiento y manipulación de datos}
\subsection{NumPy y Pandas}
% Pandas para la fusión de CSVs, limpieza de la tabla clínica y creación del Master Group ID.
% NumPy para las transformaciones matriciales de bajo nivel.
\subsection{Scikit-Learn}
% Imprescindible mencionarlo: lo usas para la normalización (StandardScaler), 
% la codificación (OneHotEncoder) y el cálculo de métricas (F1-score, Recall, AUC).

\section{Librerías de visión por computador}
\subsection{OpenCV (cv2) y Pillow (PIL)}
% Pillow para cargar las imágenes RGB estándar.
% OpenCV es CRÍTICO mencionarlo aquí porque es el motor matemático que usas 
% para generar tu rama de imágenes geométricas (Adaptive Radial Polar - ARP).

\section{Librerías de optimización e interpretabilidad}
\subsection{Optimización Bayesiana con Optuna}
% Explicar por qué usas Optuna en lugar de un GridSearch tradicional: 
% eficiencia computacional para buscar la mejor Focal Loss y Learning Rate.
\subsection{Interpretabilidad visual con pytorch-grad-cam}
% Justificar su uso para abrir la "caja negra" del modelo, verificar que no hay 
% "Target Leakage" visual (Clever Hans effect) y generar mapas de calor médicos.

\section{Entorno de ejecución y recursos hardware}
\subsection{Configuración del entorno local (VS Code)}
% Tu IDE de desarrollo.
\subsection{Limitaciones de VRAM y paralelización}
% Explicar las características de la GPU utilizada y cómo esto condicionó el 
% uso de técnicas como "Gradient Accumulation" para poder entrenar modelos 
% tan pesados (la Fusión Triple) sin quedarse sin memoria.

\section{Reproducibilidad y control de versiones}
\subsection{Fijación de semillas aleatorias (Seed locking)}
% Explicar cómo fijar la semilla en PyTorch, Numpy y Python garantiza 
% que el modelo siempre inicializa igual y los splits (5-fold) son repetibles.
\subsection{Registro de métricas y \textit{Logs}}
% Cómo se guardan los "best_model.pth" y los historiales de entrenamiento para 
% trazar las gráficas de validación.